% Part: second-order-logic
% Chapter: syntax-and-semantics
% Section: introduction

\documentclass[../../../include/open-logic-section]{subfiles}

\begin{document}

\olfileid{sol}{syn}{int}

\olsection{مقدمة}

نعبر في منطق الرتبة الأولى عن البنى بضم الرموز غير المنطقية،
أي !!{constant}s و!!{function}s و!!{predicate}s اللغة، إلى
رموزها المنطقية؛ وبذلك نصف !!{structure}s الرتبة الأولى.
وواسطة هذا الوصف مفهوم الإرضاء: يربط !!a{structure}
$\Struct{M}$ وتعيين المتغيرات $\OLId{latin-ordinary}{s}$ بـ!!a{formula} $!A$.
فمعنى $\Sat{\OLId{latin-looped}{M}}{!A}[\OLId{latin-ordinary}{s}]$ أن ما تقوله $!A$ صادق إذا أوّلت
!!{constant}s و!!{function}s و!!{predicate}s وفق $\Struct{M}$،
وأوّلت المتغيرات الحرة وفق $\OLId{latin-ordinary}{s}$. وأما الـ!!{identity} $\eq$،
فتأويلها مقرر في تعريف $\Sat{\OLId{latin-looped}{M}}{!A}[\OLId{latin-ordinary}{s}]$ نفسه، شأن
$\lforall$ و$\lexists$: الهوية هي دائمًا علاقة الهوية على
!!{domain} $\Domain{\OLId{latin-looped}{M}}$ للبنية، والمكممان يجريان على
الـ!!{domain} كله. والقيد الأساسي هنا أن التكميم لا يتناول إلا
أفراد الـ!!{domain}؛ فلا يتلو المكمم إلا !!{variable}s الأفراد.

وأما منطق الرتبة الثانية فيزيد في اللغة وفي تعريف الإرضاء
متغيرات للدوال والمحمولات، حرة ومربوطة، ويجيز التكميم عليها.
ونسبة هذه المتغيرات إلى !!{function}s و!!{predicate}s كنسبة
متغيرات الأفراد إلى !!{constant}s: تقوم مقامها في إنشاء الحدود
و!!{formula}s، ويجري تكميمها على نحو مشابه.
وفي الدلالة \emph{القياسية} يشمل مكمم الرتبة الثانية كل ما
يناسب نوع متغيره: الدوال ذات $\OLId{latin-ordinary}{n}$ مواضع من $\Domain{\OLId{latin-looped}{M}}^\OLId{latin-ordinary}{n}$
إلى $\Domain{\OLId{latin-looped}{M}}$ للمتغيرات الدالية، والعلاقات ذات $\OLId{latin-ordinary}{n}$ مواضع
للمتغيرات المحمولية. فالصيغة
$\lforall[\Obj{v_0}][(\Obj{P^1_0}(\Obj{v_0}) \lor \lnot
  \Obj{P^1_0}(\Obj{v_0}))]$ جائزة في منطقي الرتبتين؛
لكن منطق الرتبة الثانية يجيز فوقها
$\lforall[\Obj{V^1_0}][\lforall[\Obj{v_0}][(\Obj{V^1_0}(\Obj{v_0}) \lor
    \lnot \Obj{V^1_0}(\Obj{v_0}))]]$ و
$\lexists[\Obj{V^1_0}][\lforall[\Obj{v_0}][(\Obj{V^1_0}(\Obj{v_0}) \lor
    \lnot \Obj{V^1_0}(\Obj{v_0}))]]$.
وهاتان !!{sentence}s من الرتبة الثانية لخلّوهما من المتغيرات
الحرة. فـ$\Obj{V^1_0}$ متغير محمولي من الرتبة الثانية ذو $\OLMathNumeral{1}$
موضع، والقيم الجائزة لـ$\Obj{V^1_0}$ هي تأويلات
!!{predicate} ذي $\OLMathNumeral{1}$ موضع، مثل $\Obj{P^1_0}$؛ أي المجموعات
الجزئية من $\Domain{\OLId{latin-looped}{M}}$. فالتكميم يقول إن
$\lforall[\Obj{v_0}][(\Obj{V^1_0}(\Obj{v_0}) \lor \lnot
  \Obj{V^1_0}(\Obj{v_0}))]$ تصدق مهما اخترنا مجموعة جزئية من
$\Domain{\OLId{latin-looped}{M}}$ قيمةً لـ$\Obj{V^1_0}$، أو تصدق لاختيار واحد
على الأقل. ولأن الفرد إما أن ينتمي إلى المجموعة وإما ألا ينتمي،
فكلتا الجملتين صادقة في كل !!{structure}.
\end{document}
